\section{Sistem Tipe dan Pemeriksaan Tipe}
\label{sec:ch15-tipe}

\subsection{Sistem Tipe}

Sistem tipe menentukan aturan pembentukan tipe dan kompatibilitas.
Pada Subset-C, tipe dasar adalah \code{int}, \code{float},
\code{char}, \code{void}, \code{string}, dan \code{bool}.
Tipe \code{error} dipakai untuk propagasi.

\begin{figure}[H]
\centering
\begin{tikzpicture}[node distance=6mm and 10mm]
  \node[phase, text width=2.4cm] (num) {numerik};
  \node[phase, text width=2.4cm, below left=of num] (int) {int};
  \node[phase, text width=2.4cm, below=of num] (float) {float};
  \node[phase, text width=2.4cm, below right=of num] (char) {char};
  \draw[arrow] (int) -- (num);
  \draw[arrow] (float) -- (num);
  \draw[arrow] (char) -- (num);
\end{tikzpicture}
\caption{Hierarki tipe numerik Subset-C}
\label{fig:ch15-tipe}
\end{figure}

\subsection{Aturan Kompatibilitas}

Aturan assignment:
\begin{itemize}[leftmargin=*]
  \item tipe sama selalu kompatibel;
  \item \code{int} dan \code{char} dapat dipromosikan ke
        \code{float};
  \item \code{char} dapat dipromosikan ke \code{int};
  \item \code{float} tidak dapat ditetapkan ke \code{int} tanpa cast.
\end{itemize}

Aturan aritmetika:
\begin{itemize}[leftmargin=*]
  \item jika salah satu operand \code{float}, hasilnya \code{float};
  \item jika keduanya \code{int} atau \code{char}, hasilnya
        \code{int};
  \item operand non-numerik menghasilkan \code{error}.
\end{itemize}

\subsection{Traversal AST}

Pemeriksaan tipe dilakukan dengan traversal AST. Fungsi
\code{check\_expr} mengembalikan tipe simpul dan melaporkan
kesalahan. Potongannya ditunjukkan pada \lstref{lst:ch15-check-expr}.

\includecode[
  language=C,
  firstline=55,
  lastline=140,
  firstnumber=55,
  caption={Fungsi \code{check\_expr}},
  label={lst:ch15-check-expr}
]{../../code/subsetc/src/typecheck.c}

Setiap cabang menangani satu jenis simpul. Untuk identifier,
pemeriksaan dilakukan melalui tabel simbol. Untuk ekspresi biner,
tipe anak dihitung lebih dahulu, lalu aturan operator diterapkan.

\subsection{Blok dan Deklarasi}

Fungsi \code{check\_block} mendorong scope, memproses deklarasi,
lalu melepas scope. Potongannya ditunjukkan pada
\lstref{lst:ch15-check-block}.

\includecode[
  language=C,
  firstline=142,
  lastline=182,
  firstnumber=142,
  caption={Fungsi \code{check\_block}},
  label={lst:ch15-check-block}
]{../../code/subsetc/src/typecheck.c}

Deklarasi diperiksa terhadap \code{lookup\_current} sebelum
ditambahkan. Inisialisasi diperiksa terhadap tipe deklarasi. Jika
inisialisasi tidak kompatibel, kesalahan dilaporkan pada lokasi
deklarasi.

\subsection{Fungsi dan Parameter}

Fungsi didefinisikan pada scope global. Parameter didefinisikan pada
scope fungsi, lalu blok tubuh diperiksa. Tipe kembalian fungsi
disimpan pada simbol fungsi dan dipakai saat memeriksa pemanggilan.
